\documentclass[12pt]{article}
\usepackage[margin=1in]{geometry}
\usepackage{amsmath, amssymb, mathtools}
\usepackage{enumitem}
\usepackage{booktabs}
\usepackage{hyperref}
\usepackage{physics}
\hypersetup{colorlinks=true,linkcolor=black,urlcolor=blue,citecolor=blue}
\setlist[enumerate]{itemsep=0.6em}
%\setlist[enumerate,1]{label=\textbf{Q\arabic*.}}

\newcommand{\R}{\mathbb{R}}
%\newcommand{\dd}{\,\mathrm{d}}
\newcommand{\dtotal}[2]{\frac{\mathrm{d} #1}{\mathrm{d} #2}}
\newcommand{\pd}[2]{\frac{\partial #1}{\partial #2}}

\title{Solutions -- Practice Problem Set 1}
\author{ECON 320 -- Introduction to Mathematical Economics}
\date{Fall 2025}

\begin{document}
\maketitle




\begin{enumerate}
\item \textbf{Solve inequalities and express in interval notation.}
\begin{enumerate}[label=(\alph*)]
  \item $\lvert 2x-3\rvert \le 5$.\\
  By definition, $-5\le 2x-3\le 5$. Add $3$: $-2\le 2x\le 8$. Divide by $2$: $\boxed{-1\le x\le 4}$, i.e.\ $\boxed{[-1,4]}$.

  \item $(x-3)(x+2)\le 0$.\\
  Roots at $x=-2,3$. Since it is a product of two linear factors, it is $\le0$ \emph{between} the roots. Hence $\boxed{[-2,3]}$.

  \item $\dfrac{1}{x-2}\ge 3$.\\
  Case $x>2$: multiply by $x-2>0$ gives $1\ge 3(x-2)\iff 3x\le 7\iff x\le \tfrac{7}{3}$. Together with $x>2$, we get $(2,\tfrac{7}{3}]$.\\
  Case $x<2$: multiply by $x-2<0$ \emph{reverses} inequality: $1\le 3(x-2)\implies x\ge \tfrac{7}{3}$, impossible with $x<2$.\\
  Therefore $\boxed{(2,\;7/3]}$.
\end{enumerate}

\item \textbf{Set operations for $A=[-1,4)$ and $B=(1,6]$.}
\begin{enumerate}[label=(\alph*)]
  \item $A\cup B$ spans from the least left endpoint to the greatest right endpoint without gaps: $\boxed{[-1,6]}$.\\
        $A\cap B$ is the overlap: $(1,4)$. So $\boxed{A\cap B=(1,4)}$.
  \item $A\setminus B$: points in $A$ not in $B$ $\Rightarrow$ $\boxed{[-1,1]}$.\\
        $B\setminus A$: points in $B$ not in $A$ $\Rightarrow$ $\boxed{[4,6]}$.
  \item For $A\cap B=(1,4)$, $\inf=1$, $\sup=4$. Hence $\boxed{\inf=1,\ \sup=4}$.
\end{enumerate}

\item \textbf{Domain of} $\displaystyle f(x)=\frac{\sqrt{\,5-\ln(x^{2}-9)\,}}{x-3}$.\\
Conditions: (i) $\ln(x^2-9)$ defined $\Rightarrow x^2-9>0\Rightarrow |x|>3$; (ii) radicand $\ge0$ $\Rightarrow 5-\ln(x^2-9)\ge0\iff \ln(x^2-9)\le 5\iff x^2-9\le e^5\iff |x|\le \sqrt{9+e^5}$; (iii) denominator $\neq 0$ (already enforced by $x>3$ with $x\ne 3$).\\
Combine: $\boxed{\big[-\sqrt{9+e^5},\,-3\big)\ \cup\ \big(3,\,\sqrt{9+e^5}\big]}$.

\item \textbf{Inverse of} $f(x)=\dfrac{2x-3}{x+1}$ \textbf{(domain $x>-1$)}; \textbf{evaluate $f^{-1}(5)$.}\\
Solve $y=\dfrac{2x-3}{x+1}\Rightarrow y(x+1)=2x-3\Rightarrow x(y-2)=-(y+3)\Rightarrow \boxed{f^{-1}(y)=\dfrac{y+3}{2-y}}$.\\
Since $x>-1$, $f(x)=2-\dfrac{5}{x+1}$ has range $(-\infty,2)$. Therefore \textbf{$f^{-1}(5)$ is undefined} (because $5\notin(-\infty,2)$). The algebraic value $-8/3$ lies outside the domain $x>-1$.

\item \textbf{Solve for $x$.}
\begin{enumerate}[label=(\alph*)]
  \item $e^{3x}=7\Rightarrow 3x=\ln 7\Rightarrow \boxed{x=\tfrac{1}{3}\ln 7}$.
  \item $2^{x+1}=5e^{0.2x}\Rightarrow (x+1)\ln 2=\ln 5+0.2x$.\\
  $x(\ln 2-0.2)=\ln(5/2)\Rightarrow \boxed{x=\dfrac{\ln(5/2)}{\ln 2-0.2}}\approx \boxed{1.858}$.
\end{enumerate}

\item \textbf{Sequence} $S=\left\{\dfrac{2n+1}{n+3}:n\in\mathbb{N}\right\}$ (take $\mathbb{N}=\{1,2,\dots\}$).\\
Rewrite: $\dfrac{2n+1}{n+3}=2-\dfrac{5}{n+3}$. Thus $\lim_{n\to\infty}=2$. The sequence is increasing and bounded above by $2$, with smallest term at $n=1$: $3/4$. Hence $\boxed{\inf S=3/4,\ \sup S=2}$ and $\boxed{\lim S=2}$.

\item \textbf{Set} $B=\{(x,y)\in\R_{\ge0}^2:2x+3y\le 60\}$.\\
Boundary $2x+3y=60$ gives intercepts: $x=30$ (when $y=0$), $y=20$ (when $x=0$). So the feasible set is a right triangle. Area $=\tfrac12(30)(20)=\boxed{300}$. In function form: $y\le \frac{60-2x}{3}$ for $0\le x\le 30$. Hence $\boxed{B=\{(x,y):0\le x\le 30,\ 0\le y\le (60-2x)/3\}}$.

\item \textbf{Composition and transformation.} $u(x)=\ln x$, $g(x)=2-3\,\ln(4-x)$.\\
Domain: $4-x>0\Rightarrow \boxed{x<4}$. Composition: $\boxed{g(x)=2-3\,u(4-x)}$. Interpretation: start with $u(x)$, reflect horizontally across $x=2$ (since $x\mapsto 4-x$), then scale vertically by $3$ and reflect (minus sign), and shift up by $2$.

\item \textbf{Solve} $|x-a|\le b$ ($b>0$).\\
By definition, $-b\le x-a\le b\Rightarrow a-b\le x\le a+b$. Hence $\boxed{[a-b,\ a+b]}$.

\item \textbf{Domain of} $h(x)=\sqrt{9-x^2}+\dfrac{1}{\sqrt{x-1}}$.\\
Require $9-x^2\ge 0\Rightarrow -3\le x\le 3$, and $x-1>0\Rightarrow x>1$ (strict because of denominator). Intersection: $\boxed{(1,3]}$.
\end{enumerate}
\newpage


\begin{enumerate}[resume]
\item \textbf{Differentiate and evaluate.}
\begin{enumerate}[label=(\alph*)]
  \item $y=\dfrac{(3x^2+1)^5}{x e^{2x}}$. Use logarithmic differentiation:
  \[
  \ln y=5\ln(3x^2+1)-\ln x-2x\;\Rightarrow\; \frac{y'}{y}=\frac{30x}{3x^2+1}-\frac{1}{x}-2.
  \]
  At $x=1$: $y(1)=\dfrac{4^5}{e^2}=\dfrac{1024}{e^2}$ and $\frac{y'}{y}\big|_{1}=\frac{30}{4}-1-2=\frac{9}{2}$. Hence
  \[
  \boxed{y'(1)=\frac{1024}{e^2}\cdot\frac{9}{2}=\frac{4608}{e^2}}.
  \]

  \item $y=\ln\!\big(\sqrt{1+4x^2}\big)=\tfrac12\ln(1+4x^2)$. Then
  \[
  y'=\tfrac12\cdot \frac{8x}{1+4x^2}=\boxed{\frac{4x}{1+4x^2}},\qquad
  y''=\frac{(4)(1+4x^2)-4x(8x)}{(1+4x^2)^2}=\boxed{\frac{4-16x^2}{(1+4x^2)^2}}.
  \]
\end{enumerate}

\item \textbf{Cost \& average cost.} $C(q)=0.5q^2+10q+400$.\\
$MC=C'(q)=\boxed{q+10}$, $AC=C/q=\boxed{0.5q+10+\frac{400}{q}}$.\\
At $q=20$: $MC=30$, $AC=10+10+20=\boxed{40}$.\\
Minimize $AC$: $AC'(q)=0.5-400/q^2=0\Rightarrow q^2=800\Rightarrow \hat q=\boxed{20\sqrt{2}}$. Then
\[
AC(\hat q)=0.5(20\sqrt2)+10+\frac{400}{20\sqrt2}=10\sqrt2+10+10\sqrt2=\boxed{20\sqrt2+10}\ (\approx 38.284).
\]

\item \textbf{Elasticity of demand.}
\begin{enumerate}[label=(\alph*)]
  \item $Q(P)=1000P^{-1.2}$. Then $dQ/dP=-1.2\cdot 1000P^{-2.2}$. Elasticity $\varepsilon=\frac{dQ}{dP}\frac{P}{Q}=\boxed{-1.2}$ (constant).
  \item $Q(P)=200-2P$. Then $\varepsilon=\frac{-2P}{200-2P}$. At $P=50$, $Q=100$, so $\boxed{\varepsilon=-1}$ (unit elastic).
\end{enumerate}

\item \textbf{Implicit differentiation.} $x e^x+y^2=10$. Differentiate:
\[
e^x(x+1)+2y\,y'=0\ \Rightarrow\ \boxed{y'=-\dfrac{e^x(x+1)}{2y}}.
\]
At $x=1$: $e+y^2=10\Rightarrow y=\sqrt{10-e}$ (take $y>0$). Thus $\boxed{y'(1)=-\dfrac{e}{\sqrt{10-e}}}$.

\item \textbf{Differentials / percentage changes.} $Q=K^{0.3}L^{0.7}P^{0.1}$. For small changes,
\[
\frac{\dd Q}{Q}\approx 0.3\frac{\dd K}{K}+0.7\frac{\dd L}{L}+0.1\frac{\dd P}{P}.
\]
With $+5\%$, $-2\%$, $+10\%$: $\Delta Q/Q\approx 0.3(0.05)+0.7(-0.02)+0.1(0.10)=\boxed{0.011}=\boxed{+1.1\%}$.

\item \textbf{Tax $t$ in linear market.}
\[
Q_s=-20+4(P-t),\quad Q_d=700-6P.
\]
Equate: $-20+4(P-t)=700-6P\Rightarrow 10P=720+4t\Rightarrow \boxed{P(t)=72+0.4t}$. Then
\[
Q(t)=700-6P=700-6(72+0.4t)=\boxed{268-2.4t}.
\]
Hence $\boxed{\dfrac{\dd P}{\dd t}=0.4}$, $\boxed{\dfrac{\dd Q}{\dd t}=-2.4}$. Interpretation: buyers bear $40\%$ of the marginal tax (price to them rises by \$0.40 per \$1 tax).

\item \textbf{Profit maximization \& comparative statics.} $\pi(q)=pq-(aq+\tfrac{b}{2}q^2)=(p-a)q-\tfrac{b}{2}q^2$.\\
FOC: $(p-a)-bq=0\Rightarrow \boxed{q^*=\dfrac{p-a}{b}}$ (since $b>0$). Then $\boxed{\pd{q^*}{p}=1/b>0}$, $\boxed{\pd{q^*}{a}=-1/b<0}$.

\item \textbf{Elasticity of supply.} $Q(P)=cP^{\eta}\Rightarrow \dfrac{\dd Q}{\dd P}=c\eta P^{\eta-1}$. Then
\[
\varepsilon_s=\frac{\dd Q}{\dd P}\frac{P}{Q}=\frac{c\eta P^{\eta-1}P}{cP^\eta}=\boxed{\eta}.
\]
For $c=5$, $\eta=0.8$: $\boxed{\varepsilon_s=0.8}$ at any $P$ (constant).

\item \textbf{Implicit-function comparative statics.} $F(x,a)=x^3-4x+a=0$. By the IFT,
\[
\frac{\dd x}{\dd a}=-\frac{F_a}{F_x}=-\frac{1}{3x^2-4}.
\]
At $x=2$: $\boxed{\dd x/\dd a=-1/8}$.

\item \textbf{Dynamic comparative statics.}
\[
\begin{cases}
\alpha Y-\beta P=Q,\\
\gamma P-\delta W=Q.
\end{cases}
\]
Solve: $\alpha Y-\beta P=\gamma P-\delta W\Rightarrow (\beta+\gamma)P=\alpha Y+\delta W$, hence
\[
\boxed{P(Y,W)=\frac{\alpha Y+\delta W}{\beta+\gamma}},\qquad
\boxed{Q(Y,W)=\frac{\gamma\alpha}{\beta+\gamma}Y-\frac{\beta\delta}{\beta+\gamma}W}.
\]
Differentiate w.r.t.\ time:
\[
\boxed{\dtotal{P}{t}=\frac{\alpha}{\beta+\gamma}\dot Y+\frac{\delta}{\beta+\gamma}\dot W},\qquad
\boxed{\dtotal{Q}{t}=\frac{\gamma\alpha}{\beta+\gamma}\dot Y-\frac{\beta\delta}{\beta+\gamma}\dot W}.
\]
With $\alpha=\beta=\gamma=\delta=2$: $\dfrac{\dd Q}{\dd t}=\dot Y-\dot W$. Given $\dot Y=3$, $\dot W=-1$, obtain $\boxed{\dd Q/\dd t=4}$.

\item \textbf{Log-linear approximation for profits.}\\
$R(q)=Aq^\theta\Rightarrow MR=A\theta q^{\theta-1}$, $C(q)=c_0+c_1q\Rightarrow MC=c_1$. For a small change $\Delta q$ near $q_0$,
\[
\Delta\pi \approx (MR_0-MC_0)\,\Delta q=\big(A\theta q_0^{\theta-1}-c_1\big)\,(q_0\,\underbrace{\Delta q/q_0}_{\%\,\text{change}}).
\]
Numerically, $A=200$, $\theta=0.6$, $c_1=20$, $q_0=50$, $+2\%$ so $\Delta q=0.02\,q_0=1$. Compute $MR_0=200\cdot 0.6\cdot 50^{-0.4}\approx 25.095$. Thus
\[
\Delta\pi \approx (25.095-20)\cdot 1=\boxed{\approx 5.10}.
\]

\item \textbf{Subsidy $s$ in linear market.}\\
$Q_d=300-3P$, $Q_s=-30+2(P+s)=-30+2P+2s$. Equate:
\[
-30+2P+2s=300-3P\Rightarrow 5P=330-2s\Rightarrow \boxed{P(s)=66-0.4s},
\]
\[
Q(s)=300-3P=300-3(66-0.4s)=\boxed{102+1.2s}.
\]
Hence $\boxed{\dd P/\dd s=-0.4}$ (consumer price falls as subsidy rises), $\boxed{\dd Q/\dd s=1.2}$ (quantity increases).

% Q23 intentionally removed

\item \textbf{Implicit elasticity from $F(Q,P)=\ln Q+\phi Q-\ln P=0$ ($\phi>0$).}\\
Compute partials: $F_Q=1/Q+\phi$, $F_P=-1/P$. By implicit differentiation:
\[
\left(\frac{1}{Q}+\phi\right)\dd Q-\frac{1}{P}\dd P=0\ \Rightarrow\ \frac{\dd Q}{\dd P}=\frac{1/P}{\,1/Q+\phi\,}=\frac{Q/P}{1+\phi Q}.
\]
Elasticity $\displaystyle \varepsilon=\frac{\dd Q}{\dd P}\frac{P}{Q}=\boxed{\frac{1}{1+\phi Q}}$ (positive here because $F$ implies an increasing $Q$--$P$ relation).
\end{enumerate}


\end{document}
